\subsection{Degrees}

Let \(\mathbf{P} = \mathbf{A}[(\mathbf{X}_i)_{i\in I}]\) be a polynomial algebra. For each integer \(n\in \mathbf{N}\) let \(\mathbf{P}_n\) be the submodule of \(\mathbf{P}\) generated by the monomials \(\mathbf{X}^{\nu}\) such that \(|\nu| = \sum_{i\in I}\nu_i\)

equals \(n\) . Then \((\mathbf{P}_n)_{n \in \mathbf{N}}\) is a graduation which turns \(\mathbf{A}[(\mathbf{X}_i)_{i \in I}]\) into a graded algebra of type \(\mathbf{N}\) (III, p.~459). The homogeneous elements of degree \(n\) in \(\mathbf{A}[(\mathbf{X}_i)_{i \in I}]\) are sometimes called forms of degree \(n\) with respect to the indeterminates \(\mathbf{X}_i\) .

When we are dealing with the degree of inhomogeneous polynomials, we shall agree to adjoin to the set N of natural numbers, an element written \(-\infty\) and to extend the order relation and the addition of N to \(N \cup \{-\infty\}\) by the following conventions, where \(n \in N\) ,

\[
- \infty <   n, (- \infty) + n = n + (- \infty) = - \infty , (- \infty) + (- \infty) = - \infty .
\]

Let \(u = \sum_{\nu \in \mathbf{N}^{(I)}} \alpha_{\nu} X^{\nu}\) be a polynomial. The homogeneous component \(u_{n}\) of degree

\(n\) of \(u\) (for the graduation of type \(\mathbf{N}\) defined above) is equal to \(\sum_{|\nu| = n} \alpha_{\nu} X^{\nu}\) , and we

clearly have \(u = \sum_{n \in \mathbb{N}} u_n\) . If \(u \neq 0\) , the \(u_n\) are not all zero, and we define the degree

(or total degree) of u, written \(\deg u\) , as the greatest of the numbers n such that \(u_{n} \neq 0\) ; in other words (III, p.~453), the degree of u is the largest of the integers \(|\nu|\) for the multi-indices \(\nu\) such that \(\alpha_{\nu} \neq 0\) . When u = 0, the degree of u is \(-\infty\) by convention. For every integer \(p \in N\) , the relation \(\deg u \leqslant p\) is thus equivalent to \(\alpha_{\nu} = 0\) for every multi-index \(\nu\) with \(|\nu| > p\) ; the set of polynomials u such that \(\deg u \leqslant p\) is thus an A-submodule of \(\mathbf{A}[(\mathbf{X}_{i})_{i \in I}]\) , equal to \(P_{0} + P_{1} + \cdots + P_{p}\) with the above notations.

Let E be an A-module. The family \((\mathrm{E} \otimes \mathrm{P}_{n})_{n \in \mathbb{N}}\) is a graduation of type N of the module \(\mathrm{E}[(\mathrm{X}_{i})_{i \in I}] = \mathrm{E} \otimes_{\mathbb{A}} \mathrm{A}[(\mathrm{X}_{i})_{i \in I}]\) of polynomials with coefficients in E. We extend the conventions adopted above for the degree of inhomogeneous polynomials to this case.

PROPOSITION 1. --- Let u and v be two polynomials.

\begin{enumerate}
\def\labelenumi{(\roman{enumi})}
\tightlist
\item
  If \(\deg u \neq \deg v\) , we have
\end{enumerate}

\[
u + v \neq 0 \quad a n d \quad \deg (u + v) = \sup (\deg u, \deg v).
\]

If \(\deg u = \deg v\) we have \(\deg (u + v) \leqslant \deg u\) .

\begin{enumerate}
\def\labelenumi{(\roman{enumi})}
\setcounter{enumi}{1}
\tightlist
\item
  We have \(\deg(uv) \leqslant \deg u + \deg v\) .
\end{enumerate}

The proof is immediate.

Let \(J \subset I\) and \(B = A[(X_i)_{i \in I - J}]\) ; we shall identify \(A[(X_i)_{i \in I}]\) with \(B[(X_i)_{i \in J}]\) (III, p.~453 f.). The degree of \(u \in A[(X_i)_{i \in I}]\) , qua element of \(B[(X_i)_{i \in J}]\) , is called the degree of \(u\) with respect to the \(X_i\) of index \(i \in J\) (III, p.~454).

Let \(u = \sum_{k=0}^{n} \alpha_k X^k \in A[X]\) be a non-zero polynomial of degree \(n\) in a single

indeterminate. The coefficient \(\alpha_{n}\) which is \(\neq0\) by hypothesis, is called the leading coefficient of u. A polynomial \(u\neq0\) whose leading coefficient is equal to 1 is called a monic polynomial.

In \(\mathbf{A}[\mathbf{X}_1, \mathbf{X}_2, \ldots, \mathbf{X}_q]\) the number of monomials of total degree \(p\) is equal to the number of elements \((n_k)_{1 \leqslant k \leqslant q}\) of \(\mathbf{N}^q\) such that \(\sum_{k=1}^{q} n_k = p\) , that is \(\binom{q + p - 1}{p}\)

(Sets III, Prop. 15, p.~182).

More generally, let \(\Delta\) be a commutative monoid and \((\delta_i)_{i \in I}\) a family of elements of \(\Delta\) . There exists a unique graduation of type \(\Delta\) of the algebra \(\mathbf{A}[(\mathbf{X}_i)_{i \in I}]\) such that each monomial \(\mathbf{X}^\nu\) is of degree \(\sum_{i \in I} \nu_i \delta_i\) (III, p.~458,

example 3). The case considered above is that where \(\Delta = \mathbf{N}\) and \(\delta_{i} = 1\) . In the general case, to avoid confusion, we shall use the word « weight » instead of « degree » and « isobaric » instead of « homogeneous ». For example, there exists a unique graduation of type \(\mathbf{N}\) of the algebra \(\mathbf{A}[(\mathbf{X}_{i})_{i\geqslant 1}]\) such that \(\mathbf{X}_i\) is of weight \(i\) for each integer \(i \geqslant 1\) . The isobaric elements of weight \(n\) are the polynomials of the form \(\sum_{\nu} a_{\nu} X^{\nu}\) , where \(a_{\nu} = 0\) for \(\sum_{i \geqslant 1} i \cdot \nu_i \neq n\) .

